\section{Evaluation}
\label{sec:eval}

\subsection{Experimental Setup}
\label{subsec:setup}

% VALUE: one setup paragraph: the platform and its power state, the software, the workloads and prompt lengths, the baseline, how TTFT is measured, the correctness gate, and how the machine state is controlled; defaults stated once
We measure on a GMKtec NucBox EVO-X2 with an AMD Ryzen\texttrademark{} AI Max+ 395 processor, whose XDNA~2 NPU has 8 columns of 4 compute tiles, and 64\,GB of shared system memory.
The host runs Linux~7.0 with the amdxdna driver and XRT~2.25.0~\citep{xdna-driver,xrt-repo}, the NPU in its performance power mode, and the CPU's frequency governor and energy preference set to performance.
We evaluate Llama-3.2-1B~\citep{meta2024llama32} and Qwen3-4B~\citep{yang2025qwen3}, and \E{a third, 8B-class model}, with 4-bit quantized weights in FastFlowLM's format~\citep{fastflowlm-repo}.
Prompts computed from scratch have 2{,}048 or 4{,}096 tokens, and appends add \E{512, 1{,}024,} 2{,}048, or 4{,}096 tokens to a cached context of 4{,}096, 16{,}384, or 28{,}672 tokens, followed by 250 generated tokens in an agent's turn.
These appends span the median to beyond the 90th percentile of the new tokens per step in coding-agent traces~\citep{zhu2026tracelab}, whereas the contexts stop at the 32{,}768-token window we build, below the traces' median prefix of 119K tokens, so our results cover agents whose context fits an edge model's window.
% SOURCE: TraceLab Table 8: append tokens P25 / P50 / P90 = 384 / 857 / 5,342 (Claude) and 406 / 886 / 8,310 (Codex); output P50 252 / 184; Section 1: median prefix 119K, append 875, output 214 (PDF read 2026-10-15).
% DECISION: the 512 and 1,024-token appends enter once measured (queued); until then they stay \E.
The designs are MLIR-AIR's fused prefill and fused decode~\citep{wang2026mlirair}, compiled with MLIR-AIE~\citep{mlir-aie-repo} and used unchanged.
The baseline is each design's own host program on its default split, which runs the same configuration, kernels, and weights as our composed schedules.
Every run's logits must be bit-identical to the baseline's, every append's logits bit-identical to those of a full prefill of the whole prompt, and the kernel log must show no driver or IOMMU message during the run.
TTFT is the wall time of the prefill call, from token identifiers to logits, and we report the median of \E{3} timed runs after one warm run, with the order of schedules rotated between rounds.
Decode time per token is measured over 64 generated tokens at a given context length.
Each batch of runs holds the device exclusively, checks that no other process has it open, and brackets the runs with a DRAM bandwidth probe and an empty-command dispatch.
Because the aggregate DRAM read cap varies by up to 12\% between batches, every ratio compares runs within one batch.
% SOURCE: machine files read 2026-10-14: uname -r 7.0.0-34-generic; amdxdna 2.25.0_20260922 (DKMS, 31c01f2019e7); XRT 2.25.0 (3923b2b5be55); BIOS EVO-X2 1.12; NPU power mode performance, CPU EPP and governor performance (lab-prefill-overlap (11)).
% SOURCE: mlir-air a3e26261 (llms/shared/fused_prefill); mlir_aie wheel 1.4.5.dev22+g35fa6b1; weights FastFlowLM/Qwen3-4B-NPU2 rev 07651f11cac7 and the Llama-3.2-1B q4nx bundle, both sha256-verified (lab-prefill-overlap (23)).
% SOURCE: runs per point today: 3 for the from-scratch and Llama-3.2-1B append comparisons, 2 for the split comparison and the Qwen3-4B appends (vault lab-prefill-overlap (18)-(26), lab-agent-step (2), lab-split-sweep (4)-(5)); decode, 64 timed tokens (lab-decode-long-context (5)).
% DECISION: the run count is fixed at the results freeze; a count that differs by comparison is stated where it applies.
% SOURCE: lock = flock ~/.local/state/npu.lock; brackets = 1x1 and 8x2 DRAM reads of 256 MiB, the dispatch floor, load, NPU mode, EPP, governor; 8x2 cap 49.9-55.7 GB/s over the batches, 55.7 / 49.9 = 1.116 (lab-prefill-device (1), lab-prefill-overlap (2), (11), lab-split-sweep (4)-(5), lab-agent-step (2)).
% DECISION: the NPU firmware version and the BIOS power mode are not yet recorded; add them to the setup once the author records them.

\subsection{Instruction Order}
\label{subsec:ttft}
% VALUE: composing the instruction order lowers TTFT on every model and length, more where attention's share is larger, with bit-identical outputs; on appends to a cached context the gain holds at a steady 1.5-1.6x; the existing order's slowdown at positions ~8k-20k on Llama-3.2-1B is explained, not quoted as a gain; the essential decisions are reported apart from host-boundary handling (queue depth, packing); the gap to the lower bound per configuration
% DECISION: results stay \E until the results freeze (2026-10-21) and the audit's verdicts.
\E{TTFT from scratch and on appends, by model and length; the steps from the existing schedule through co-scheduled groups, predicted-completion order, and host-boundary handling; the gap to the lower bound.}

\subsection{Spatial Split}
\label{subsec:split-eval}
% VALUE: the faster split changes with the regime on both models (two attention columns from scratch, four for appends to a long context), and the gain of choosing both the split and the order grows with the cached context; switching between resident splits costs a suspend and resume, under 1% of an append; the four-column split's gain comes from attention despite slower GEMMs
\E{TTFT per split, schedule, and regime on both models; the switch cost.}

\subsection{Model Accuracy}
\label{subsec:accuracy}
% VALUE: the event model prices schedules and a model it was not fitted on within a few percent, picks the faster split in every configuration measured, and states where its error grows (attention read off its trend, the existing order's mid-range slowdown) with the size of each
\E{Prediction error on schedules and models not used to build the model, per component; the split choice; the regime where the error grows.}

\subsection{Generality Across Models}
\label{subsec:generality}
% VALUE: the same composer and rules apply unchanged to a second family and an 8B-class model, with bit-identical outputs; the arena rule splits the larger weights
\E{Second and third models.}

\subsection{Agent Turns}
\label{subsec:turns}
% VALUE: measured whole turns (an append plus its decode) on both models; decode's share of the turn, characterized in Section 3.2, bounds the turn-level gain
% DECISION: turn-level numbers are stated only once measured; until then this subsection stays a placeholder.
\E{Measured whole turns on both models.}
